\BourbakiExerciseGroup

\begin{enumerate}
\def\labelenumi{\arabic{enumi})}
\item
  设 E 是域 K 的一个扩张，B 是 E 在 K 上的一个超越基。证明：若集合 K、B 中有一个是无限的，则 E 与 \(K \times B\) 等势，否则 E 是可数的。* 特别地，由此推出实数域 R 在有理数域 Q 上的每个超越基都具有连续统的势。*
\item
  设 K 为 \(\mathbf{Q}(\mathbf{X})\) ，即有理数域 Q 上一个未定元的有理分式域。证明：在环 K{[}Y{]} 中，多项式 \(Y^{2} + X^{2} + 1\) 不可约；又若 E 是由该多项式的一个根（在
\end{enumerate}

K），则 Q 在 E 中相对代数封闭，但 E 不是 Q 的纯扩张。（要看出 E 中没有在 Q 上代数的元素 \(a \notin Q\) ，注意这样一个元素的存在将蕴含关系 \(\mathrm{E} = \mathbf{Q}(a)(\mathbf{X})\) ，并注意 \(-(\mathbf{X}^{2} + 1)\) 在 \(\Omega(\mathbf{X})\) 中不是平方元，其中 \(\Omega\) 是 Q 的一个代数闭包。）证明：若 i 是 \(X^{2} + 1\) 的一个根，则 \(\mathrm{E}(i)\) 是 \(\mathbf{Q}(i)\) 的纯超越扩张。

* 3) 设 K 为域 C(X)，即（代数闭的）复数域 C 上一个未定元的有理分式域。证明：在环 K{[}Y{]} 中，多项式 \(Y^{3} + X^{3} + 1\) 不可约；设 E 为由该多项式的一个根（在 K 的一个代数闭包中）生成的 K 的扩张。证明 E 不是 C 的纯扩张，尽管 C 是代数闭的。（证明：关系 \(u^{3} + v^{3} + w^{3} = 0\) 不可能存在于 C{[}X{]} 的三个两两互素且并非全为常数的多项式 u、v、w 之间。用反证法论证：若 r 是 u、v、w 的次数中的最大者，且例如 \(\deg(w) = r\) ，则由关系式

\[
w ^ {3} = - (u + v) (u + j v) (u + j ^ {2} v)
\]

其中 \(j\) 是本原三次单位根，由该关系式推出存在三个多项式 \(u_1\) 、\(v_1\) 、\(w_1\) （在 \(\mathbf{C}[\mathbf{X}]\) 中），它们两两互素、并非全为常数、次数 \(< r\) ，且满足 \(u_1^3 + v_1^3 + w_1^3 = 0\) 。）*

\begin{enumerate}
\def\labelenumi{\arabic{enumi})}
\setcounter{enumi}{3}
\tightlist
\item
  设 \(\Omega\) 是域 \(\mathbf{K}\) 的一个代数闭扩张，且在 \(\mathbf{K}\) 上具有无限超越次数。
\end{enumerate}

\begin{enumerate}
\def\labelenumi{\alph{enumi})}
\tightlist
\item
  证明：存在无穷多个 \(\Omega\) 的 K-自同态，其像是 \(\Omega\) 的异于 \(\Omega\) 的子域，并且关于这些自同态，\(\Omega\) 可以具有等于直至 tr. \(\deg_{\mathbf{K}}\Omega\) 的任意基数的超越次数。* 特别地，存在无穷多个不同的复数域 C 到 C 的异于 C 的子域上的 Q-同构。* b) 证明：对 \(\Omega\) 的每个满足 tr. \(\deg_{\mathbf{K}}E < \text{tr} \cdot \deg_{\mathbf{K}}\Omega\) 的子扩张 E，E 到 \(\Omega\) 的一个子域上的每个 K-同构都可以延拓为 \(\Omega\) 的一个 K-自同构。给出这样的例子：其中 tr. \(\deg_{\mathbf{K}}E = \text{tr} \cdot \deg_{\mathbf{K}}\Omega\) ，且存在 E 到 \(\Omega\) 的一个子域上的一个 K-同构，它不能延拓为 E 的任一扩张 F（包含于 \(\Omega\) 且异于 E）到 \(\Omega\) 的一个子域上的 K-同构。
\end{enumerate}

\begin{enumerate}
\def\labelenumi{\arabic{enumi})}
\setcounter{enumi}{4}
\item
  设 \(\Omega\) 是域 \(\mathbf{K}\) 的一个代数闭扩张。证明：对每个子扩张 \(\mathbf{E}\) —— \(\Omega\) 的、在 \(\mathbf{K}\) 上为超越的 —— 存在无穷多个 \(\mathbf{K}\) -同构，把 \(\mathbf{E}\) 映到 \(\Omega\) 的一个子域上（若 \(x \in \mathbf{E}\) 在 \(\mathbf{K}\) 上是超越的，考虑子扩张 \(\mathbf{F}\) （含于 \(\mathbf{E}\) ），使得 \(x\) 在 \(\mathbf{F}\) 上超越且 \(\mathbf{E}\) 在 \(\mathbf{F}(x)\) 上代数，并证明存在无穷多个 \(\mathbf{F}\) -同构，从 \(\mathbf{E}\) 到 \(\Omega\) 的一个子域上）。
\item
  设 \(\Omega\) 是域 \(K\) 的一个代数闭扩张，\(N\) 是 \(\Omega\) 的一个子扩张。证明：若 \(N\) 在 \(K\) 上是拟伽罗瓦的，且 \(E\) 与 \(E'\) 是 \(K\) 的两个包含在 \(\Omega\) 中的共轭扩张，则 \(N(E)\) 与 \(N(E')\) 在 \(K\) 上共轭。反之，若 \(N\) 具有此性质，且 \(\operatorname{tr} \cdot \deg_{K} N\) 有限并 \(< \operatorname{tr} \cdot \deg_{K} \Omega\) ，则它必然在 \(K\) 上是代数的且是拟伽罗瓦的。
\item
  设 E 是域 K 的一个扩张，\(\operatorname{Aut}_{\mathbf{K}}(\mathbf{E})\) 是 E 的所有 K-自同构组成的群。a) 对 E 在 K 上的每个有限生成子扩张 L，\(\operatorname{Aut}_{\mathbf{L}}(\mathbf{E})\) 是 \(\operatorname{Aut}_{\mathbf{K}}(\mathbf{E})\) 的一个子群。证明：当 L 遍历 E 的在 K 上有限生成的子扩张所成的集合时，诸群 \(\operatorname{Aut}_{\mathbf{L}}(\mathbf{E})\) 构成 \(\mathrm{Aut}_{\mathbf{K}}(\mathbf{E})\) 的单位元的一个基本邻域系，对应于一个与 \(\mathrm{Aut}_{\mathbf{K}}(\mathbf{E})\) 的群结构相容的拓扑，且对该拓扑 \(\mathrm{Aut}_{\mathbf{K}}(\mathbf{E})\) 是一个分离且完全不连通的群。
\end{enumerate}

\begin{enumerate}
\def\labelenumi{\alph{enumi})}
\setcounter{enumi}{1}
\tightlist
\item
  当 E 赋有离散拓扑时，\(\operatorname{Aut}_{\mathbf{K}}(\mathbf{E})\) 的拓扑是使映射 \((u,x)\to u(x)\) —— 从 \(\operatorname{Aut}_{\mathbf{K}}(\mathbf{E})\times\mathbf{E}\) 到 E 的映射 —— 连续的最粗拓扑。
\end{enumerate}

* c) 群 \(\operatorname{Aut}_{\mathbf{Q}}(\mathbf{R})\) 约化为中性元素（V，第 153 页，§9 习题 2）。* \(d)\) 设 \(\mathbf{K}_0 \supseteq \mathbf{K}\) 是 E 的由在 E 的所有 K-自同构下不变的元素组成的子扩张。证明：要使 E 在 \(\mathbf{K}_0\) 上是代数的，必须而且只需 \(\operatorname{Aut}_{\mathbf{K}}(\mathbf{E})\) 是紧致的。

\begin{enumerate}
\def\labelenumi{\alph{enumi})}
\setcounter{enumi}{4}
\tightlist
\item
  假设 K 是无限的。若 \(E = K(X)\) 是 K 上一个未定元的有理分式域，证明 \(\operatorname{Aut}_{K}(E)\) 是一个无限离散群（V，第 148 页，习题 11）。若 K 的特征为 0，则 \(\Gamma\) —— \(\operatorname{Aut}_{K}(E)\) 中由自同构 \(\sigma : X \mapsto X + 1\) 生成的子群 —— 在 \(\operatorname{Aut}_{K}(E)\) 中是闭的且异于后者，但具有相同的不变域。
\end{enumerate}

\begin{enumerate}
\def\labelenumi{\arabic{enumi})}
\setcounter{enumi}{7}
\tightlist
\item
  设 \(E\) 是域 \(K\) 的一个扩张。
\end{enumerate}

\begin{enumerate}
\def\labelenumi{\alph{enumi})}
\item
  设 \(x\) 是 \(E - K\) 的一个元素。证明：存在一个纯扩张 \(L \subset E\) （ \(K\) 的），使得 \(E\) 在 \(L\) 上是代数的，且 \(x \notin L\) （考虑子扩张 \(L\) —— 属于 \(E\) 、在 \(K\) 上纯超越且 \(x \notin L\) 的 —— 所成的集合）。
\item
  设 \(x \in E\) 是 \(K\) 上的一个超越元，\(p\) 是一个整数 \(\geqslant 2\) 。证明：存在一个纯超越扩张 \(L \subset E\) （ \(K\) 的），使得 \(E\) 在 \(L\) 上是代数的，\(x \notin L\) 且 \(x^p \in L\) （同样的方法）。
\item
  设 F 是 E 的一个子扩张，\(x \in E\) 是 F 上的一个代数元，\(P \in F[X]\) 是它的极小多项式。证明存在 K 的一个纯扩张 \(L \subset E\) ，使得 E 在 L 上是代数的，且 P 在域 \(F(L)\) 上不可约。
\end{enumerate}

\begin{enumerate}
\def\labelenumi{\arabic{enumi})}
\setcounter{enumi}{8}
\tightlist
\item
  域 K 的一个扩张 E 称为 K 的戴德金扩张，如果对 E 的每个子扩张 F，F 都与 E 中在群 \(\mathrm{Aut}_{\mathrm{F}}(\mathrm{E})\) —— E 的所有 F-自同构组成的群 —— 作用下不变的元素所构成的子域相同。
\end{enumerate}

\begin{enumerate}
\def\labelenumi{\alph{enumi})}
\item
  要使 \(\mathbf{E}\) —— \(\mathbf{K}\) 的代数扩张 —— 在 \(\mathbf{K}\) 上是戴德金的，必须而且只需 \(\mathbf{E}\) 是 \(\mathbf{K}\) 的一个伽罗瓦扩张。
\item
  假设 K 的特征为 p \textgreater{} 0。证明 K 的戴德金扩张必然在 K 上是代数的（从而是伽罗瓦的）。（利用习题 8，b）。）
\item
  如果E是K的戴德金扩张，且L是E的一个纯超越子扩张，不同于K，并且E是L的代数扩张，证明E在L上的次数是无限的。
\end{enumerate}

\begin{enumerate}
\def\labelenumi{\arabic{enumi})}
\setcounter{enumi}{9}
\item
  设 E 是域 K 的一个扩张，使得映射 \(F \mapsto \operatorname{Aut}_{F}(E)\) 是从 E 的子扩张所成的集合到 \(\operatorname{Aut}_{K}(E)\) 的子群所成的集合上的一个双射。证明此时 E 是 K 的一个有限次伽罗瓦扩张。（利用习题 9，c)，先证明 E 在 K 上必然是代数的，再利用习题 9，a)。）
\item
  设 \(\mathbf{E}\) 是域 \(\mathbf{K}\) 的一个扩张，\(\Omega\) 是 \(\mathbf{E}\) 的一个代数闭扩张。证明下列条件是等价的：
\end{enumerate}

α) E是K的p-根式扩张；

\(\beta)\) 对每个扩张 \(F\) （在 \(K\) 之上），\(E \otimes_{K} F\) 只有唯一的素理想；

γ) 对于K的每个扩张F，E与F的任意两个复合扩张都是同构的；

δ) \(E \otimes_{K} \Omega\) 只有一个素理想；

ε) E 和 Ω 的任意两个复合扩张都是同构的。

\begin{enumerate}
\def\labelenumi{\arabic{enumi})}
\setcounter{enumi}{11}
\tightlist
\item
  \begin{enumerate}
  \def\labelenumii{\alph{enumii})}
  \tightlist
  \item
    设 K 是一个域，\(\Omega\) 是 K 的一个代数闭扩张，\(E \subset \Omega\) 与 F 是 K 的两个扩张；假设 \(\operatorname{tr} \cdot \deg_E \Omega \geqslant \operatorname{tr} \cdot \deg_K F\) 。证明 E 与 F 的每个复合扩张都同构于形如 \((L_w, 1, w)\) 的复合，其中 \(w\) 是 F 到 \(\Omega\) 的一个子域上的一个 K-同构，而 \(L_w\) 是 \(\Omega\) 中由 E 与 \(w(F)\) 生成的子域。
  \end{enumerate}
\end{enumerate}

\begin{enumerate}
\def\labelenumi{\alph{enumi})}
\setcounter{enumi}{1}
\tightlist
\item
  由 a) 推出：若 F 在 K 上是 n 次有限代数扩张，则 E 与 F 的复合扩张中两两不同构的至多有 n 个。
\end{enumerate}

\begin{enumerate}
\def\labelenumi{\arabic{enumi})}
\setcounter{enumi}{12}
\item
  在 \(\Omega\) —— \(\mathbf{Q}(\mathbf{X})\) 的一个代数闭包 —— 中，考虑 \(\mathrm{E} = \mathbf{Q}(\mathbf{X})\) 与 \(\mathrm{F} = \mathbf{Q}(\mathbf{X} + i)\) （其中 \(i^2 = -1\) ）这两个域 \(\mathbf{Q}\) 的纯超越扩张。证明 \(\mathrm{E} \cap \mathrm{F} = \mathbf{Q}\) ，但 \(\mathrm{E}\) 与 \(\mathrm{F}\) 在 \(\mathbf{Q}\) 上并非代数不相交。
\item
  设 E、F、G 是域 K 的包含于 K 的一个扩张 \(\Omega\) 中的三个扩张，且满足 \(F \subset G\) 。证明：要使 E 与 G 在 K 上代数不相交，必须而且只需 E 与 F 在 K 上代数不相交，并且 \(\mathrm{E}(\mathrm{F})\) 与 G 在 F 上代数不相交。
\item
  \begin{enumerate}
  \def\labelenumii{\alph{enumii})}
  \tightlist
  \item
    设 K 是一个域，L 是 K 的一个子域，使得 K 是一个有限生成的 L-代数，换言之，\(K = L[a_{1}, a_{2}, \ldots, a_{n}]\) ，其中元素属于 K。证明这些 \(a_{j}\) 在 L 上是代数的。（用反证法论证：若 \(a_{1}, \ldots, a_{m}\) （ \(m \geqslant 1\) ）构成 \((a_{j})_{1 \leqslant j \leqslant n}\) 的一个极大代数无关子族，注意在环 \(A = L[a_{1}, \ldots, a_{m}]\) 中所有非零理想的交集为零，并利用 V，第 147 页，习题 5。）
  \end{enumerate}
\end{enumerate}

\begin{enumerate}
\def\labelenumi{\alph{enumi})}
\setcounter{enumi}{1}
\tightlist
\item
  由 a) 推出：在多项式代数 \(K[X_{1}, \ldots, X_{n}]\) 中，每个极大理想在 K 上的余维数都是有限的。由此得出结论：若 A 是 K 上的一个交换有限生成代数，则 A 的包含 K 的每个子域在 K 上的次数都是有限的。
\end{enumerate}

\begin{enumerate}
\def\labelenumi{\arabic{enumi})}
\setcounter{enumi}{15}
\tightlist
\item
  \begin{enumerate}
  \def\labelenumii{\alph{enumii})}
  \tightlist
  \item
    设 K 是一个代数闭域，\(\Omega\) 是 K 的一个代数闭扩张，E 是 \(\Omega\) 的一个子扩张，\(x\) 是 \(\Omega\) 的一个在 E 上超越的元素。设 \(\overline{\mathbf{K}(x)}\) 是 \(\mathbf{K}(x)\) 在 \(\Omega\) 中的相对代数闭包，并令 \(\mathbf{M} = \operatorname{E}(\overline{\mathbf{K}(x)})\) ；证明 M 是 \(\operatorname{E}(x)\) 的一个无限次代数扩张。（注意，对每个整数 \(m > 1\) ，多项式 \(\mathbf{X}^m - x\) 在 \(\operatorname{E}(x)[\mathbf{X}]\) 中不可约（V，第 91 页，注记）。
  \end{enumerate}
\end{enumerate}

\begin{enumerate}
\def\labelenumi{\alph{enumi})}
\setcounter{enumi}{1}
\tightlist
\item
  设 E 是一个代数闭域，F 是 E 的一个有限生成扩张。证明 F 的每个代数闭子域必然包含在 E 中。（否则将存在一个代数闭的子域 \(L \subset F\) ，以及一个在 E 上超越的元素 \(x \in L\) 。应用 a)，取 K 为 E 的素子域在 E 中的相对代数闭包，并注意 M 应在 E 上有限生成。）
\end{enumerate}

\begin{enumerate}
\def\labelenumi{\arabic{enumi})}
\setcounter{enumi}{16}
\item
  设 K 是一个域，\(\Omega\) 是 K 的一个代数闭扩张，x 是 \(\Omega\) 的一个在 K 上超越的元素。证明：对每个素数 q，存在 \(\Omega\) 的一个子扩张 M，使得 \(\mathbf{M}(x)\) 是 M 的 q 次伽罗瓦扩张。（若 q 不是 K 的特征，考虑元素 \(y = x^{q}\) ，在相反的情形考虑 \(y = x^{q} - x\) ；再考虑 \(\Omega\) 的一个子扩张 L，它由 \(\Omega\) 在 K 上的一个包含 y 的超越基生成，且若 \(y = x^{q}\) 则还加上 q 次单位根；利用 V，第 163 页，习题 12。）
\item
  设 K 是一个域，\(\mathbf{K}((\mathbf{X}))\) 是 K 上一个未定元的形式幂级数域（IV，第 38 页）。证明 \(\operatorname{tr}.\deg_{\mathbf{K}}\mathbf{K}((\mathbf{X}))=\operatorname{Card}(\mathbf{K}^{\mathbf{N}})\) 。区分两种情形：
\end{enumerate}

\begin{enumerate}
\def\labelenumi{\alph{enumi})}
\item
  若 \(\operatorname{Card}(\mathbf{K}) < \operatorname{Card}(\mathbf{K}^{\mathbf{N}})\) ，注意 \(\operatorname{Card}(\mathbf{K}((\mathbf{X}))) = \operatorname{Card}(\mathbf{K}^{\mathbf{N}})\) 。
\item
  若 \(\operatorname{Card}(K) = \operatorname{Card}(K^{N})\) ，设 P 是 K 的素子域，S 是 \(\operatorname{K}((X))\) 中一个在 K 上代数无关的无限元素集，T 是属于 S 的所有形式幂级数的系数所成的集合。设 L 是域 K(S) 在
\end{enumerate}

K((X)) 中的相对代数闭包，并设 u 是 L 的一个元素；存在一个方程 \(g(s_{1}, \ldots, s_{m}, u) = 0\) ，其中 \(g \neq 0\) 是 \(K[X_{1}, \ldots, X_{m}, X_{m+1}]\) 中的一个多项式，而 \(s_{1}, \ldots, s_{m}\) 是 S 的元素。设 A 是多项式 g 的系数所成的集合，\(C(u)\) 是形式幂级数 u 的系数所成的集合；证明域 \(P(T)(A \cup C(u))\) 在 \(P(T \cup A)\) 上是代数的（用 \(\Omega\) 表示 K 的一个代数闭包，证明否则将存在无穷多个形式幂级数 \(v \in \Omega((X))\) 满足方程 g（ \(s_{1}, \ldots, s_{m}, v) = 0\) ）。利用习题 1，证明若 \(\text{Card}(S) < \text{Card}(K)\) ，则 K 在 \(P(T)\) 上的超越次数是无限的，并由此推出在这种情形 L 异于 \(K((X))\) （若 \((t_{n})_{n \geqslant 0}\) 是 K 中在 \(P(T)\) 上代数无关的一个无穷序列，考虑形式幂级数 \(\sum_{n=0}^{\infty} t_{n} X^{n}\) ）。

\begin{enumerate}
\def\labelenumi{\arabic{enumi})}
\setcounter{enumi}{18}
\tightlist
\item
  设 \(\mathbf{E}\) 与 \(\mathbf{F}\) 是域 \(\mathbf{K}\) 的两个在 \(\mathbf{K}\) 上线性不相交的超越扩张。证明 \(\mathbf{K}(\mathbf{E} \cup \mathbf{F})\) 异于环 \(\mathbf{C}\) （同构于 \(\mathbf{E} \otimes_{\mathbf{K}} \mathbf{F}\) ，由 \(\mathbf{E} \cup \mathbf{F}\) 生成的）。（归结为 \(\mathbf{E}\) 与 \(\mathbf{F}\) 在 \(\mathbf{K}\) 上的超越次数都等于 1 的情形；若 \(x \in \mathbf{E}\) 与 \(y \in \mathbf{F}\) 在 \(\mathbf{K}\) 上是超越的，证明我们不可能有 \((x + y)^{-1} \in \mathbf{C}\) ；用反证法论证，证明在相反的情形，若 \(p\) 是 \(\mathbf{K}\) 的特征指数，则将存在整数 \(r \geqslant 0\) 使得 \((x + y)^{-p^r}\) 属于 \(\mathbf{C}\) 的由 \(\mathbf{K}(x) \cup \mathbf{K}(y)\) 生成的子环。）
\end{enumerate}

* 20) 设 \(n\) 是一个整数，\(K\) 是一个特征 \(p\) 与 \(n\) 互素、且具有本原 \(n\) 次单位根的域，\(\mu_n \subset K^*\) 是 \(n\) 次单位根所成的子群，\(G\) 是一个被 \(n\) 零化的有限交换群，而 \(A = K^{(G)}\) 是 \(G\) 的群代数（III，第 446 页）。

\begin{enumerate}
\def\labelenumi{\alph{enumi})}
\item
  证明 A 是一个可对角化代数（V，第 28 页）。（将 G 分解为循环群的直和（VII，第 22 页）。）由此推出：每个 A-模都是一些模的直和，这些模是维数为 1 的向量 K-空间。
\item
  设 V 是 K 上有限维的 A-模，S 是向量 K-空间 V 的对称代数，F 是 S 的分式域。群 G 作用在 V 上，从而作用在 S 与 F 上。假设 G 在 V 上的作用是忠实的。证明 G 在 F 上的作用是忠实的。计算 \([F : F^{G}]\) 。
\item
  设 \((x_{1},\ldots ,x_{n})\) 是向量 K-空间 V 的一组由 G 的作用的特征向量组成的基。设 \(\Gamma\) 是 \(\mathbf{F}^*\) 的由 \(x_{1},\ldots ,x_{n}\) 生成的子群。证明 \(\Gamma\) 是一个自由交换群，且 F 的由 \(\Gamma\) 生成的子 K-代数 B 可与 \(\mathbf{K}^{(\Gamma)}\) 等同，并且在 G 的作用下是稳定的。
\item
  对所有 \(x \in \Gamma\) ，映射 \(\gamma \mapsto \frac{\gamma x}{x}\) 是从 \(G\) 到 \(\mu_n\) 中的一个同态，由此给出一个同态 \(\theta: \Gamma \to \operatorname{Hom}(G, \mu_n)\) 。证明 \(\theta\) 是满射的。令 \(\Gamma_1 = \operatorname{Ker}(\theta)\) 。利用初等因子定理（VII，第 24 页），证明 \(B\) 是一个秩为 Card(G) 的自由 \(K^{(\Gamma_1)}\) -模。
\item
  通过过渡到分式域，由 \(d\) 与 \(a\) 推出 \(\mathbf{F}^{\mathrm{G}}\) 是 \(\mathbf{K}^{(\Gamma_1)}\) 的分式域。由此推出 \(\mathbf{F}^{\mathrm{G}}\) 是 \(\mathbf{K}\) 的一个纯超越扩张。*
\end{enumerate}

